% !TEX ROOT = ../Thesis.tex 
%-
%-> 中文摘要
%-
\chapter{摘\quad 要}\chaptermark{摘\quad 要}% 摘要标题
\pagenumbering{Roman}% 页码符号
\setcounter{page}{2}% 开始页码
% 22p / 12p = 1.83
\linespread{1.5}
\zihao{-4}
% \setlength{\baselineskip}{20pt}

在“双碳”战略持续推进和交通运输体系加速电气化的背景下，电动汽车进入规模化应用阶段。动力电池与旋转机械作为支撑整车动力输出的关键部件，在长期复杂工况下易发生性能退化与结构故障，直接影响车辆运行安全、系统可靠性与服役寿命。依托多源传感数据构建数据驱动的智能预测与健康管理方法，已成为保障电动汽车全生命周期高可靠运行的重要技术路径。然而，受观测噪声、工况变化、域分布漂移和未知输入等因素影响，现有深度模型在实际部署中普遍存在推断偏差、分布外泛化能力不足以及大语言模型幻觉等问题，难以满足安全关键场景下的高可信决策需求。针对上述问题，本文以无偏不确定性量化为主线，面向电动汽车关键部件开展退化建模与故障诊断研究，构建覆盖可信状态估计、分布外鲁棒退化建模与抗幻觉故障诊断的智能预测与健康管理方法体系，
% 
本文的具体贡献如下：

（1）针对锂离子电池健康状态可信估计问题，提出了一种基于层次化随机时空 Transformer 的健康状态估计方法。
首次，该方法利用 Gumbel-Softmax 重参数化构建随机自注意力机制，使特征聚合过程具备不确定性量化能力。
其次，结合计及不确定性的逐时间步卷积、距离敏感的层次化 Transformer 编码器，实现时空不确定性的量化和时空退化特征的提取。
最后，利用基于保序回归的不确定性校准策略，实现时空退化特征的有效提取与预测置信度的无偏修正。
实验结果表明，所提出方法在多个锂离子电池循环老化数据集上均取得优于对比方法的估计精度与不确定性量化效果，并能够有效辨识分布内与分布外退化模式。

（2）针对锂离子电池在未知工况与未见退化模式下的泛化退化建模问题，提出了一种基于拓扑约束下不确定性驱动特征调弈的分布外鲁棒退化建模方法。
首先，该方法通过启发式循环信号重构与基于 Mamba 的流形收缩特征提取策略，构建低本征维度且具有拓扑连续性的退化表征。
其次，结合多元高斯统计建模与不确定性驱动特征合成策略，对潜在分布外退化模式进行刻画，
最后，设计长尾感知的自适应重平衡优化机制，以缓解样本分布不均引发的建模偏置。
基于多个开源电池数据集构建的分布外泛化实验表明，所提出方法在退化建模任务上表现出更强的鲁棒性与泛化能力。

（3）针对多模态大语言模型在旋转机械故障诊断中的幻觉问题，提出了一种基于扩散增强可信多模态大语言模型的抗幻觉故障诊断方法。
首先，该方法构建 ControlNet 引导的跨数据集条件扩散模型，生成多样化少样本故障信号，以缓解类别统计偏差。
其次，设计文本与信号对齐的可信多模态大语言模型，并通过关键 token 定位与不确定性重加权策略实现序列级幻觉检测。
最后，构建面向非预期输入的抗幻觉诊断流程，提升模型对未知输入的拒识能力与可靠诊断能力。
实验结果表明，所提出方法在多个故障诊断数据集和多项评价指标上均优于现有对比方法。

本文围绕电动汽车关键部件在复杂工况下的高可靠运行需求，系统研究无偏不确定性量化在可信估计、分布外泛化与抗幻觉诊断中的作用机理与实现方法。相关研究不仅提升数据驱动模型的预测精度与可靠性，也为构建面向实际应用场景的电动汽车智能预测与健康管理系统提供理论依据与方法支撑。

\keywords{电动汽车; 预测与健康管理; 深度学习; 锂离子电池; 旋转机械}% 中文关键词

%-
%-> 英文摘要
%-
\chapter[Abstract]{{\rmfamily Abstract}}\chaptermark{Abstract}% 摘要标题

With the continued advancement of the carbon peaking and carbon neutrality strategy and the accelerated electrification of transportation, electric vehicles have entered large-scale deployment. Power batteries and rotating machinery, as key components of the vehicle powertrain, are prone to degradation and faults under long-term complex operating conditions, directly affecting operational safety, reliability, and service life. Data-driven prognostics and health management has therefore become an important technical pathway for reliable electric vehicle operation. However, observation noise, operating-condition variations, distribution shifts, and unknown inputs often cause inference bias, weak out-of-distribution generalization, and hallucinations in large language models, making existing methods inadequate for safety-critical deployment. To address these issues, this dissertation focuses on unbiased uncertainty quantification and investigates degradation modeling and fault diagnosis for key electric vehicle components, forming a framework that covers trustworthy state estimation, out-of-distribution robust degradation modeling, and hallucination-resistant fault diagnosis.

The main contributions of this dissertation are as follows:

(1) For trustworthy state-of-health estimation of lithium-ion batteries, a hierarchical stochastic spatiotemporal Transformer is proposed. The method introduces Gumbel-Softmax-based stochastic self-attention for uncertainty-aware feature aggregation, combines uncertainty-aware step-wise convolution with a distance-aware hierarchical Transformer encoder, and uses isotonic-regression-based calibration to correct predictive confidence. Experiments on multiple battery aging datasets show superior estimation accuracy and uncertainty quantification performance, as well as effective identification of in-distribution and out-of-distribution degradation patterns.

(2) For generalized lithium-ion battery degradation modeling under unknown operating conditions and unseen degradation patterns, a topology-constrained uncertainty-driven feature adaptation method is proposed. It constructs degradation representations with low intrinsic dimensionality and topological continuity through heuristic cycling signal reconstruction and Mamba-based manifold-shrinking feature extraction, characterizes potential out-of-distribution patterns with multivariate Gaussian modeling and uncertainty-driven feature synthesis, and alleviates modeling bias through long-tail-aware adaptive rebalancing. Experiments on multiple open-source battery datasets demonstrate stronger robustness and generalization capability.

(3) For hallucinations in multimodal large language models for rotating machinery fault diagnosis, a diffusion-enhanced trustworthy multimodal large language model is proposed. A ControlNet-guided cross-dataset conditional diffusion model generates diverse few-shot fault signals to reduce category-level statistical bias. Text-signal alignment, key-token localization, and uncertainty reweighting are then used for sequence-level hallucination detection, while an anti-hallucination diagnostic pipeline improves rejection and reliability under unexpected inputs. Experiments show that the proposed method outperforms existing methods on multiple fault diagnosis datasets and metrics.

In summary, this dissertation investigates unbiased uncertainty quantification for trustworthy estimation, out-of-distribution generalization, and hallucination-resistant diagnosis of key electric vehicle components. The studies improve the accuracy and reliability of data-driven models and provide methodological support for practical intelligent prognostics and health management systems.

\englishkeywords{Electric vehicles; Prognostics and health management; Deep learning; Lithium-ion batteries; Rotating machinery; Large language models}
